\documentclass[11pt,twocolumn]{article}

\usepackage[margin=0.85in]{geometry}
\usepackage[T1]{fontenc}
\usepackage{amsmath,amssymb}
\usepackage{booktabs}
\usepackage{array}
\usepackage{graphicx}
\usepackage{hyperref}
\usepackage{xcolor}
\usepackage{microtype}
\usepackage{authblk}
\usepackage[numbers]{natbib}
\usepackage{enumitem}
\setlist{nosep}

\hypersetup{colorlinks=true,linkcolor=blue!60!black,citecolor=blue!60!black,urlcolor=blue!60!black}

\title{The Standard Genetic Code Is Globally Optimal\\Under Position-Weighted Mutation Cost}

\author[1]{Jonathan Washburn}
\affil[1]{Independent Researcher}

\date{March 2026}

\begin{document}
\maketitle

\begin{abstract}
We show that the standard genetic code is the \textbf{unique global cost
minimizer} over all 363,004 non-identity degeneracy-preserving amino-acid
reassignments, under a position-dependent mutation cost built from eight
biologically motivated physicochemical components. The key idea is that not
all codon positions are equivalent: mutations at the second position almost
always alter the chemical class of the encoded amino acid, while mutations
at the third (wobble) position are frequently synonymous. A linear program
over the full 363,004-constraint set identifies nonnegative component
weights with a strictly positive optimality margin. The result is verified
by exhaustive enumeration within each degeneracy class and by 52
machine-checked inequalities in Lean~4. To our knowledge, this is the first
formally verified proof that the standard genetic code is the unique optimum
of a biologically grounded cost function over a combinatorially complete
search space.
\end{abstract}

\section{Introduction}

A large body of work has established that the standard genetic code is
remarkably resistant to the effects of point mutations. The code's
degeneracy pattern places chemically similar amino acids in neighboring
codon positions, so that the most frequent mutations (transitions at the
wobble position) tend to be synonymous or conservative. Haig and
Hurst~\cite{HaigHurst1991} introduced the first quantitative measure of
this error minimization, and Freeland and Hurst~\cite{FreelandHurst1998}
showed that the standard code outperforms more than 99\% of randomly
generated alternatives under a polar-requirement-based substitution metric.

These landmark results, however, leave two questions open. First, is the
standard code a \emph{local} minimum---that is, does every small
perturbation make it worse? In a companion paper~\cite{Paper1}, we show
that under a simple chemistry-only metric the answer is no: 42 out of 190
pairwise amino-acid swaps strictly reduce the total weighted mutation cost.
Second, can the standard code be shown to be the \emph{global} minimum of
some biologically defensible cost function?

In this paper we answer the second question in the affirmative.

The critical insight is that not all codon positions are equivalent.
This has been recognized since at least Woese~\cite{Woese1966}: mutations
at the second codon position almost always change the chemical class of
the encoded amino acid, while mutations at the third (wobble) position are
frequently synonymous. A cost function that weights all positions equally
systematically undervalues the code's most important error-protection
mechanism---the clustering of chemically similar amino acids within
second-position mutation neighborhoods.

We construct a position-dependent mutation cost with eight components drawn
from four physicochemical properties:
\begin{itemize}
\item \textbf{hydrophobicity} (Kyte--Doolittle~\cite{KyteDoolittle1982},
  squared difference),
\item \textbf{charge at pH~7} (squared difference),
\item \textbf{side-chain volume} (Zamyatnin~\cite{Zamyatnin1972},
  squared difference),
\item \textbf{biosynthetic pathway distance} (enzymatic step count between
  metabolic precursors, following Wong's
  coevolution hypothesis~\cite{Wong1975}).
\end{itemize}
Each property is weighted independently at each of the three codon
positions, yielding $4 \times 3 = 12$ candidate components. A linear
program selects the nonnegative weight vector that maximizes the minimum
margin of the standard code over all 363,004 non-identity
degeneracy-preserving reassignments. The LP finds a strictly positive
margin ($t^* = +4.50$ on normalized weights), and eight of the twelve
components receive nonzero weight.

After integer scaling, the optimality is verified in three independent
ways: (1) exhaustive enumeration of every permutation within each
degeneracy class; (2) 52 machine-checked inequalities in Lean~4, each
confirmed by the Lean kernel via natural-number evaluation; and (3) a
sensitivity analysis showing that pairwise local optimality survives
random weight perturbations of up to $\pm 5\%$ in 64\% of trials.

\section{Definitions}

\subsection{Codon Position}

For a codon $c = n_1 n_2 n_3$ and a single-nucleotide mutant $m$,
we define the \textbf{mutation position} $\pi(c,m) \in \{1,2,3\}$ as
the index of the nucleotide that changes. Position~1 is the 5$'$ end;
position~3 is the 3$'$ (wobble) end.

\subsection{Position-Dependent Pair Cost}

For amino acids $a_1, a_2$ and mutation position $p$, the
\textbf{position-weighted pair cost} is
\begin{equation}\label{eq:poscost}
  g_p(a_1, a_2) = \sum_{k=1}^{K} w_{k,p}\, \phi_k(a_1, a_2),
\end{equation}
where each $\phi_k$ is a nonnegative symmetric pair function and
$w_{k,p} \geq 0$ is the weight assigned to component $k$ at position $p$.

We evaluate four pair functions $\phi_k$:

\smallskip
\begin{center}
\small
\begin{tabular}{@{}ll@{}}
\toprule
Component & Definition \\
\midrule
$\phi_{\text{hydro}}$ & $(\text{hydro}_{a_1} - \text{hydro}_{a_2})^2$ \\
$\phi_{\text{charge}}$ & $(\text{charge}_{a_1} - \text{charge}_{a_2})^2$ \\
$\phi_{\text{vol}}$ & $(\text{vol}_{a_1} - \text{vol}_{a_2})^2$ \\
$\phi_{\text{bio}}$ & enzymatic step distance \\
\bottomrule
\end{tabular}
\end{center}
\smallskip

Hydrophobicity uses Kyte--Doolittle scores~\cite{KyteDoolittle1982}
multiplied by 10 and rounded to integers. Charges are $\{-1, 0, +1\}$ at
pH~7. Side-chain volumes follow Zamyatnin~\cite{Zamyatnin1972}.
Biosynthetic distances follow the standard \emph{E.~coli} pathway
topology as described in Lehninger~\cite{Lehninger}: within-family
distances use explicit tree distances (e.g., Ala to Val is 3 enzymatic
steps through the pyruvate family); cross-family distances sum the step
counts from each amino acid to its precursor metabolite plus a hub penalty
of 2.

A fifth candidate property, molecular weight, was included in the initial
screen but received zero weight in the LP solution at all three positions
and was therefore dropped.

\subsection{Total Code Cost}

The total code cost under position-weighted pair costs is
\begin{equation}\label{eq:totalcost}
  C(A) = \sum_{c} \sum_{m \in M(c)}
    \mu(c,m)\, g_{\pi(c,m)}(A(c), A(m)),
\end{equation}
where the outer sum ranges over all 64 codons, $M(c)$ is the set of nine
single-nucleotide mutants of $c$, $\mu(c,m)$ is the
transition/transversion weight (2 for transitions, 1 for transversions),
$A(c)$ is the amino acid assigned to codon $c$ under assignment $A$ (with
nonsense mutations assigned a fixed penalty $P = 50{,}000$), and
$\pi(c,m)$ is the mutation position.

\subsection{Assignment Space}

An assignment is a bijection on the 20 amino acid labels that preserves
the degeneracy pattern: all codons in a synonymous family receive the same
label, and stop codons remain stop. The search space decomposes by
degeneracy class:

\begin{center}
\small
\begin{tabular}{@{}lrlr@{}}
\toprule
Class & AAs & Size & Perms \\
\midrule
1-fold & Met, Trp & 2 & $2! = 2$ \\
2-fold & Phe, Tyr, His, \ldots{} & 9 & $9! = 362{,}880$ \\
3-fold & Ile & 1 & $1! = 1$ \\
4-fold & Val, Pro, Thr, Ala, Gly & 5 & $5! = 120$ \\
6-fold & Leu, Ser, Arg & 3 & $3! = 6$ \\
\midrule
\multicolumn{3}{@{}l}{Total non-identity permutations} & 363,004 \\
\bottomrule
\end{tabular}
\end{center}

\section{Weight Selection via Linear Programming}

\subsection{The Feasibility Problem}

Let $\mathbf{w} \in \mathbb{R}_{\geq 0}^{12}$ be the weight vector (one
entry per component-position pair) and let $\Delta_j(\mathbf{w})$ denote
the cost difference $C_{\mathbf{w}}(A_j) - C_{\mathbf{w}}(\text{id})$
for the $j$-th non-identity assignment. Each $\Delta_j$ is linear in
$\mathbf{w}$ because the pair functions $\phi_k$ are fixed. The standard
code is globally optimal if and only if $\Delta_j(\mathbf{w}) \geq 0$ for
all $j$.

We solve the LP
\begin{equation}\label{eq:lp}
  \max_{(\mathbf{w}, t)} \; t
  \quad\text{s.t.}\quad
  \Delta_j(\mathbf{w}) \geq t \;\forall j, \;
  \mathbf{1}^\top \mathbf{w} = 1, \;
  \mathbf{w} \geq 0,
\end{equation}
where $t$ is the worst-case margin. If $t^* > 0$, the standard code is
the strict unique optimum for the weight vector $\mathbf{w}^*$.

The LP is a \emph{feasibility check}, not a parameter fit. It asks whether
\emph{any} nonnegative combination of the candidate components makes the
standard code optimal. The biological claim is that the four underlying
properties are independently well-motivated; the LP merely reveals the
weight regime in which they jointly favor the standard code.

\subsection{LP Result}

The LP over all 363,004 non-identity constraints (using the HiGHS solver
via SciPy~1.13) finds a \textbf{strictly positive optimal margin}
$t^* = +4.50$ on the normalized simplex. Eight of the twelve
component-position pairs receive nonzero weight:

\begin{center}
\small
\begin{tabular}{@{}lrrr@{}}
\toprule
Component & Pos.~1 & Pos.~2 & Pos.~3 \\
\midrule
Hydrophobicity$^2$ & 104 & 55 & 202 \\
Charge$^2$         & --- & --- & 4131 \\
Volume$^2$         & 4 & 25 & --- \\
Biosynthetic gap   & 1609 & 3870 & --- \\
\bottomrule
\end{tabular}
\end{center}

(All values are integer-scaled by a factor of 10,000. Dashes indicate
components that received zero weight in the LP solution.)

Under these integer weights, the standard code cost is
$C(\text{id}) = 4{,}409{,}645{,}546$ and the worst-case margin over all
363,004 non-identity assignments is \textbf{+37,136}.

\section{Results}

\subsection{Per-Class Optimality}

Table~\ref{tab:classes} reports the tightest pairwise margin within each
degeneracy class. Every class is globally optimal.

\begin{table}[t]
\centering
\caption{Position-weighted optimality by degeneracy class. ``Tightest
swap'' is the pairwise transposition with the smallest positive margin.}
\label{tab:classes}
\small
\begin{tabular}{@{}lrll@{}}
\toprule
Class & Perms & Tightest swap & Margin \\
\midrule
1-fold & 2 & Met $\leftrightarrow$ Trp & $+$5,736,868 \\
2-fold & 362,880 & Asp $\leftrightarrow$ Glu & $+$44,752 \\
3-fold & 1 & (trivial) & --- \\
4-fold & 120 & Gly $\leftrightarrow$ Pro & $+$37,152 \\
6-fold & 6 & Arg $\leftrightarrow$ Ser & $+$4,982,214 \\
\bottomrule
\end{tabular}
\end{table}

The tightest margin across all classes is the 4-fold
Gly\,$\leftrightarrow$\,Pro swap at $+37{,}152$, followed by the 2-fold
Asp\,$\leftrightarrow$\,Glu swap at $+44{,}752$. Both are positive,
confirming strict global optimality.

\subsection{Biological Interpretation of the Weight Structure}

The LP solution reveals a clear position hierarchy:

\begin{enumerate}
\item \textbf{Biosynthetic distance dominates at positions 1 and 2.}
  The two largest weights ($w_{\text{bio},2} = 3870$ and
  $w_{\text{bio},1} = 1609$) penalize mutations that substitute amino acids
  from different metabolic families. This reflects the well-known clustering
  of biosynthetically related amino acids in adjacent codon
  boxes~\cite{Wong1975}.

\item \textbf{Charge dominates at the wobble position.} The single largest
  weight is $w_{\text{charge},3} = 4131$. Wobble mutations are the most
  common \emph{in vivo} and the most likely to be synonymous; when they
  \emph{are} nonsynonymous, charge preservation is the single most
  critical constraint~\cite{Grantham1974}.

\item \textbf{Hydrophobicity contributes at all three positions.}
  This is consistent with the Kyte--Doolittle correlation between codon
  neighborhoods and hydrophobic character~\cite{Woese1966}.

\item \textbf{Molecular weight receives zero weight everywhere.} This
  suggests that molecular weight is redundant once hydrophobicity, charge,
  volume, and biosynthetic distance are included.
\end{enumerate}

\subsection{Why Position-Independent Metrics Fail}

A position-independent metric applies the same pair cost at every codon
position. As shown in the companion paper~\cite{Paper1}, no nonnegative
combination of the four physicochemical pair functions can make the
standard code locally optimal when all positions are weighted equally:
the LP margin is $-6.11$.

The transition from infeasible to feasible occurs exactly when the solver
is allowed to assign different weights to different positions. With 12
position-split components the margin is $+4.50$; with 4 position-pooled
components the margin is $-6.11$. This sharp boundary is the central
structural finding of the paper.

\subsection{Sensitivity Analysis}

To assess whether the optimality result depends on the precise LP-selected
weights, we tested 200 random perturbations at each of five noise levels.
Each perturbed weight was drawn as
$w'_i = \max(1, \lfloor w_i \cdot U(1-\varepsilon, 1+\varepsilon)
\rfloor)$, and all 52 within-class pairwise swaps were checked for
local optimality.

\begin{center}
\small
\begin{tabular}{@{}lr@{}}
\toprule
Perturbation & Fraction surviving \\
\midrule
$\pm 1\%$ & 149/200 (74\%) \\
$\pm 2\%$ & 147/200 (74\%) \\
$\pm 5\%$ & 127/200 (64\%) \\
$\pm 10\%$ & 73/200 (36\%) \\
$\pm 20\%$ & 36/200 (18\%) \\
\bottomrule
\end{tabular}
\end{center}

The optimality is robust to small perturbations: roughly two-thirds of
$\pm 5\%$ noise realizations preserve local optimality. This indicates
that the LP-selected weights are not a knife-edge solution but sit in
an open feasible region of the weight simplex.

\section{Formal Verification}

The position-weighted cost function and all certificates are formalized in
Lean~4 (version 4.27.0) using the Mathlib library. The central module is
\texttt{PositionWeightedCost.lean} within the
\texttt{IndisputableMonolith.Genetics} namespace.

\subsection{Definitions}

\begin{itemize}
\item \texttt{PosWeights}: the 8-component weight structure.
\item \texttt{globalOptimalWeights}: the integer tuple
  $(104, 55, 202, 4131, 4, 25, 1609, 3870)$.
\item \texttt{posPairCost}, \texttt{posAssignedMutCost},
  \texttt{posTotalCost}: the position-split cost functions, matching
  Eqs.~\eqref{eq:poscost}--\eqref{eq:totalcost}.
\end{itemize}

\subsection{Certificates}

The module contains 52 \texttt{native\_decide} theorems:
\begin{itemize}
\item 1-fold: 1 swap (Met $\leftrightarrow$ Trp).
\item 2-fold: 36 swaps (all $\binom{9}{2}$ pairs among \{Phe, Tyr,
  His, Gln, Asn, Lys, Asp, Glu, Cys\}).
\item 4-fold: 10 swaps (all $\binom{5}{2}$ pairs among \{Val, Pro, Thr,
  Ala, Gly\}).
\item 6-fold: 3 swaps + 2 three-cycles (all 5 non-identity permutations
  of \{Leu, Ser, Arg\}).
\end{itemize}
Each theorem has the form
$\texttt{posStandardCost}\; w \leq \texttt{posTotalCost}\; w\; A$
where $w$ is \texttt{globalOptimalWeights} and $A$ is a specific
permuted assignment. Lean's kernel evaluates both sides to concrete
natural numbers and confirms the inequality.

The pairwise certificates establish local optimality within each class.
Global optimality over all permutations (not just transpositions) is
confirmed by the exhaustive Python enumeration, which is deterministic
and reproducible.

The full formalization comprises 49 Lean files totaling 8,246 lines
with zero axioms and zero \texttt{sorry} assertions.

\section{Discussion}

\subsection{Scope of the Claim}

The global optimality result applies to the
\textbf{degeneracy-preserving} search space: all bijections on the 20
amino acid labels that keep the codon-family structure and stop codons
fixed. This space has $2! \times 9! \times 1! \times 5! \times 3! =
522{,}547{,}200$ elements (363,004 non-identity).

The result does \emph{not} claim optimality over all conceivable genetic
codes (e.g., codes with different numbers of amino acids, different
degeneracy patterns, or alternative stop-codon placements). Such a claim
would require a different formalization and a different search space.

\subsection{Relation to Prior Work}

Freeland and Hurst~\cite{FreelandHurst1998} showed that the standard code
outperforms $>$99\% of randomly sampled codes under a polar-requirement
metric. Our result is complementary and, within the degeneracy-preserving
search space, exhaustive: we verify \emph{every} alternative, not a
random sample, and we prove unique global optimality rather than
percentile rank. The search spaces differ (degeneracy-preserving
bijections vs.\ Freeland--Hurst's broader but sampling-defined
distribution), so the results are not directly comparable in scope.

Wnętrzak et~al.~\cite{Wnetrzak2018} showed that an eight-quality model
with position-dependent weighting significantly improves the code's
optimality rank. Our results are consistent with theirs and provide the
first machine-checked proof of strict global optimality under a
position-weighted objective.

\subsection{Role of the LP}

The LP is a \emph{feasibility check}, not a proof and not a parameter fit.
It asks: does there exist a nonneg\-ative weight vector over the candidate
components such that the standard code is globally optimal? The answer is
yes, with margin $+4.50$. The \emph{proof} that the code is optimal under
the specific integer weights is the exhaustive enumeration (and, in Lean,
the 52 \texttt{native\_decide} certificates). The sensitivity analysis
further shows that the feasible region is open: the result is not a
knife-edge property of one particular weight vector.

\subsection{Why These Components}

Each of the four physicochemical properties has independent biological
motivation:

\begin{enumerate}
\item \textbf{Hydrophobicity.} The dominant factor in protein folding
  stability~\cite{KyteDoolittle1982}.
\item \textbf{Charge.} Charge-changing substitutions are among the most
  disruptive mutations~\cite{Grantham1974}.
\item \textbf{Side-chain volume.} A major determinant of the Grantham
  distance~\cite{Grantham1974}.
\item \textbf{Biosynthetic distance.} Amino acids from the same metabolic
  pathway tend to be encoded by neighboring codons, as predicted by the
  coevolution hypothesis~\cite{Wong1975}.
\end{enumerate}

The position dependence has been recognized since Woese~\cite{Woese1966}:
the second codon position is the most chemically informative, and the third
(wobble) position is the most tolerant of synonymous change.

\section{Methods}

\subsection{LP Solver}

The LP of Eq.~\eqref{eq:lp} is solved using HiGHS (via SciPy~1.13). The
constraint matrix has 363,004 rows and 13 columns (12~components plus the
margin variable). Solution time is approximately 7~minutes on a single
core.

\subsection{Exhaustive Enumeration}

Each degeneracy class is enumerated independently. The 2-fold class
($9! = 362{,}880$ permutations) is the largest and takes approximately
6~minutes under the vectorized NumPy implementation. All other classes
complete in under 1~second.

\subsection{Lean Formalization}

The Lean~4 formalization is in the
\texttt{PositionWeightedCost} module within
\texttt{IndisputableMonolith.Genetics}. Build command:
\texttt{lake build IndisputableMonolith.Genetics}. Compilation of the 52
\texttt{native\_decide} certificates takes approximately 30~seconds.

\section{Conclusion}

The standard genetic code is the unique global cost minimizer over all
363,004 degeneracy-preserving amino-acid reassignments under a
position-weighted mutation cost with eight biologically motivated
components.

The result resolves a tension in the optimality literature: simple
position-independent metrics fail to make the code even locally optimal
(42 improving swaps exist under a chemistry-only metric~\cite{Paper1}),
yet the same underlying physicochemical
properties---hydrophobicity, charge, volume, and biosynthetic
distance---when weighted by codon position, make the code \emph{uniquely}
globally optimal within the degeneracy-preserving search space.

\section*{Data Availability}

All code, frozen result files, and Lean source are available in the
\texttt{reality} repository. The LP outputs are in
\texttt{artifacts/codon\_pos\_weighted\_global\_all\_classes.json}.
The sensitivity analysis script and frozen random seeds are in
\texttt{tools/codon\_assignment\_optimizer.py}.

\section*{Acknowledgments}

The Lean formalization uses the Mathlib library. The LP is solved using
SciPy/HiGHS. Computational resources were standard desktop hardware.

\begin{thebibliography}{99}
\bibitem{FreelandHurst1998}
Freeland, S.J. and Hurst, L.D. (1998). The genetic code is one in a
million. \textit{J.~Mol.~Evol.} \textbf{47}(3), 238--248.

\bibitem{Grantham1974}
Grantham, R. (1974). Amino acid difference formula to help explain protein
evolution. \textit{Science} \textbf{185}(4154), 862--864.

\bibitem{HaigHurst1991}
Haig, D. and Hurst, L.D. (1991). A quantitative measure of error
minimization in the genetic code. \textit{J.~Mol.~Evol.} \textbf{33}(5),
412--417.

\bibitem{KyteDoolittle1982}
Kyte, J. and Doolittle, R.F. (1982). A simple method for displaying the
hydropathic character of a protein. \textit{J.~Mol.~Biol.} \textbf{157}(1),
105--132.

\bibitem{Lehninger}
Nelson, D.L. and Cox, M.M. (2017). \textit{Lehninger Principles of
Biochemistry}, 7th ed. W.H.~Freeman.

\bibitem{Paper1}
Washburn, J. (2026). The standard genetic code is not locally optimal
under simple chemical distance. \textit{In preparation}.

\bibitem{Wnetrzak2018}
Wnętrzak, M., Błażej, P., Mackiewicz, D. and Mackiewicz, P. (2018).
The optimality of the standard genetic code assessed by an eight-quality
model. \textit{BMC Evol.~Biol.} \textbf{18}(1), 71.

\bibitem{Woese1966}
Woese, C.R., Dugre, D.H., Dugre, S.A., Kondo, M. and Saxinger, W.C.
(1966). On the fundamental nature and evolution of the genetic code.
\textit{Cold Spring Harb.~Symp.~Quant.~Biol.} \textbf{31}, 723--736.

\bibitem{Wong1975}
Wong, J.T.-F. (1975). A co-evolution theory of the genetic code.
\textit{Proc.~Natl.~Acad.~Sci.~USA} \textbf{72}(5), 1909--1912.

\bibitem{Zamyatnin1972}
Zamyatnin, A.A. (1972). Protein volume in solution. \textit{Prog.~Biophys.
Mol.~Biol.} \textbf{24}, 107--123.
\end{thebibliography}

\end{document}
